\section{Related Work}
\label{sec:related}

\paragraph{Transformer reproduction.}
The original Transformer~\cite{vaswani2017attention} reports tokenized BLEU on WMT14 en-de (Base 27.3, Big 28.4) and en-fr (Big 41.0); subsequent re-implementations using sacrebleu~\cite{post2018call} typically land 1.5--2 BLEU lower in equivalent units, attributable to tokenizer and BLEU-variant differences~\cite{post2018call}. Our numbers (Base en-fr 35.31, Big en-fr 35.87, Base en-de 24.04 sacrebleu) are within or above this expected band.

\paragraph{Scaling laws.}
\cite{kaplan2020scaling} and \cite{hoffmann2022chinchilla} establish that test loss scales as a power law in compute, parameters, and data. These laws assume data quality is fixed across the scaling range; our 2$\times$2 directly probes the quality--capacity interaction that the laws abstract away. The closest scaling-law observation we are aware of is~\cite{tay2022scaling}'s finding that data quality interacts with effective compute, but their setting is encoder-only and does not isolate the capacity sign-flip we report.

\paragraph{Data curation for MT.}
Aggressive corpus filtering for MT is well studied~\cite{koehn2018findings, junczys2019cscope}, often using length, language-ID, and parallelism heuristics, and increasingly using neural QE filters~\cite{batheja2022qe}. We use a strict v1 filter inspired by~\cite{gururangan2020dont}'s strict-DAPT setup and a CometKiwi-22~\cite{rei2022cometkiwi} top-K filter for SFT. Our negative SFT result is consistent with~\cite{junczys2019cscope}'s finding that domain-adaptive filtering can degrade out-of-domain BLEU even when filter-quality is verifiable.

\paragraph{Loss-based sample selection.}
RHO-Loss~\cite{mindermann2022prioritized} uses model-relative loss to identify learnable samples; our spike-guard mechanism is mechanically simpler (drop high-loss batches via a 1.3$\times$EMA threshold) and is intended for noise filtering, not curriculum design. \S\ref{sec:methodology} reports an unanticipated property of this kind of mechanism: trigger rates are a joint function of data noise and model convergence depth, complicating their use as a pure data diagnostic.
